\section{Preliminaries}
\label{sec:prelims}
\xhdr{Graph search} Suppose that we are given an unweighted connected graph $\mathcal{G} = (\VS, \mathcal{E})$, where $\VS$ is a set of nodes, and $\mathcal{E}$ a corresponding set of edges. Further suppose that each node $i \in \VS$ has corresponding features $x_i \in \mathbb{R}^{D_v}$, and each edge $(i, j) \in \mathcal{E}$ has features $e_{ij} \in \mathbb{R}^{D_e}$. Assume that we are also given a start node $v_s \in \VS$ and a goal node $v_g \in \VS$. At any stage of our search algorithm, we can partition the nodes of our graph into three sets as $\VS = \text{CLOSE} \cup \text{OPEN} \cup \text{REST}$, where CLOSE are the nodes already explored,  $\text{OPEN}$ are candidate nodes for exploration (i.e., all nodes connected to any node in CLOSE, but not yet in CLOSE), and $\text{REST}$ is the rest of the graph. Each \textit{expansion} moves a node from $\text{OPEN}$ to CLOSE, and adds the neighbors of the given node from $\text{REST}$ to $\text{OPEN}$. We call the set of newly added fringe nodes $\VS_{new}$ at each search step.  At the start of the search procedure, $\text{CLOSE} = \{v_s\}$ and we expand the nodes until $v_g$ is encountered (i.e., until $v_g \in \text{CLOSE}$).

\xhdr{Greedy best-first search} We can perform \textit{greedy best-first search} using a greedy fringe expansion policy, such that we always expand the node $v \in \text{OPEN}$ that minimizes $h(v, v_g)$. Here, $h : \VS \times \VS \longrightarrow \mathbb{R}$ is a tailored heuristic function for a given use case. In our work, we are interested in learning a function $h$ that predicts shortest path lengths, this way minimizing $|\text{CLOSE}|$ in a \textit{greedy best-first search} regime.

%\update{Best-first should be clearer now - also organized this section a bit more}


%jure{We never define what a clairvoyant oracle is. Isn't oracle already optimal/supernatural/} \update{We decide to just use oracle for symbol brevity and it's easier to understand.}
\xhdr{Imitation of perfect heuristics} Partially observable Markov decision processes (POMDPs) are a suitable framework to describe the problem of learning search heuristics~\cite{bhardwaj2017learning}. We can have $s = (\text{CLOSE}, \text{OPEN}, \text{REST})$ as our state, an action $a \in \mathcal{A}$ corresponds to moving a node from $\text{OPEN}$ to CLOSE, and the observations $o \in \mathcal{O}$ are the features of newly included nodes in $\text{OPEN}$. Note that one could consider an MDP framework to learn heuristics, but the time complexity of operating on the whole state is in most cases prohibitive. We also define a history $\psi \in \Psi$ as a sequence of observations $\psi = o_1, o_2, o_3, ...$. Our work leverages the observation that using a heuristic function during greedy best-first search that correctly determines the length of the shortest path between fringe nodes and the goal node will also yield minimal $|\text{CLOSE}|$. For training, we adopt a perfect heuristic $h^{*}$, similar to \cite{bhardwaj2017learning}, which has full information about $s$ during search. Such oracle can provide ground-truth distances $h^{*}(s, v, v_g)$, where $v \in \text{OPEN}$. To conclude, we define a \textit{greedy best-first search policy} $\pi_{\theta}$ that uses a parameterized heuristic $h_{\theta}$ to expand nodes from $\text{OPEN}$ with minimal heuristic values. One could also directly use a POMDP solver for the above-described problem, but this approach is usually infeasible due to the dimensionality of the search state \cite{choudhury2017adaptive}.

%\update{Policy defined}

%Consider a partially observable Markov decision process (POMDP) defined by an 8-tuple $(S, A, R, P, \Omega, O, \gamma, \rho_0)$. In this POMDP definition, we have $S$ being a set of valid states, $A$ set of valid actions, $R : S \times A \longrightarrow \mathbb{R}$ the reward function, $P : S \times A \longrightarrow \mathcal{P}(S)$ the transition probability function, $\rho_0$ the initial state distribution, $\gamma \in [0, 1]$ the discount factor, $\Omega$ the set of observations, and $O : S \times A \longrightarrow \mathcal{P}(\Omega)$ a conditional observation probability function. We can also define a history $\psi \in \Psi$ as a sequence of observations $\psi = o_1, o_2, o_3, ...$. In general, our goal is to find an action policy $\pi(\cdot | \psi)$ that maximizes expected discounted rewards under the induced trajectory distribution. 
